% ===================================================================
%  TAULUKKO N:o 12 — Asema. — Rautatiet. — Alukset.
%  Traduction 2026 d'apres texto/io, controlee sur texto/fr et
%  texto/sv. Consigne : voir texto/fi/10-taulukko-01.tex.
% ===================================================================

%%K t12-apar-1 apar
\VUsaut{4.0mm}
\VUtitre{15.0pt}{60}{TAULUKKO N:o 12}
\VUsaut{2.4mm}
\VUpk{11.4pt}{40}{Asema. --- Rautatiet. --- Alukset.}
\VUsaut{3.4mm}
\textsc{Huom.} --- \textit{Kussakin maassa käytettävät aikataulut ovat erilaisia,
käytämme siis esimerkkinä} \VUgras{ranskalaisen taulukon} \textit{tunteja.}

%%K t12-01-1 p
1. --- Olen juuri matkustanut Saksan halki. Ennen lähtöä ostin
\VUgras{aikataulukirjan} \textsuperscript{(15)} tietääkseni junien lähtöajan.
Todettuani, että sopivin juna lähtisi Pariisista kello yhdeksän illalla ja
saapuisi Strasbourgiin kello yksi ja kolmetoista minuuttia, valmistelin matkani,
ja seuraavana päivänä annoin viedä itseni asemalle.

%%K t12-02-1 p
2. --- Kun menin suureen lähtöhalliin vanhan \VUgras{maalaispapin}
\textsuperscript{(2)} takana, joka kantoi \VUgras{matkalaukkuaan}
\textsuperscript{(3)} kädessään, monta \VUgras{matkustajaa}
\textsuperscript{(1)} oli jo tullut.

%%K t12-02-2 p
Koska halusin lähettää \VUgras{sähkeen} \textsuperscript{(5)}, annoin
\VUgras{tarkastajan} \textsuperscript{(6)} näyttää minulle
\VUgras{lennätinkonttorin} \textsuperscript{(4)}.

%%K t12-03-1 p
3. --- Ennen kuin menin \VUgras{odotussaliin} \textsuperscript{(7)}, menin
hankkimaan itselleni meno-paluu\VUgras{lipun} \textsuperscript{(9)}.

%%K t12-04-1 p
4. --- En odottanut kauan \VUgras{luukulla} \textsuperscript{(8)}, sillä siellä
oli vähän ihmisiä. \VUgras{Sotilas} \textsuperscript{(10)}, joka matkusti
lomalle, oli niin kohtelias, että jätti minulle paikkansa jonossa, niin että
minulla oli riittävästi aikaa pyytää muutamia tietoja
\VUgras{asemapäälliköltä} \textsuperscript{(13)}, joka keskusteli valvonnan
\VUgras{komisarion} \textsuperscript{(12)} ja päivystävän \VUgras{santarmin}
\textsuperscript{(11)} kanssa.

%%K t12-tit-1 sub
\VUpk{10.2pt}{40}{Matkatavaroiden kirjaaminen.}
\VUsaut{3.0mm}

%%K t12-05-1 p
5. --- Ostettuani sanomalehden \VUgras{kirjakaupasta} \textsuperscript{(14)},
annoin punnita \VUgras{matkatavarani} \textsuperscript{(17)}, jotka
\VUgras{kantaja} \textsuperscript{(16)} oli juuri tuonut sinne
\VUgras{tavarakärryllä} \textsuperscript{(19)}; \VUgras{virkailija}
\textsuperscript{(22)}, joka vastaa \VUgras{vaa'asta} \textsuperscript{(21)},
ilmoitti minulle, että matka-arkkuni painaa kolmekymmentäviisi kilogrammaa; siitä
tuli viiden kilogramman \VUgras{ylipaino}, josta olen velkaa lisämaksun,
matkatavaroiden \VUgras{luukulla} \textsuperscript{(23)}.

%%K t12-06-1 p
6. --- Koska olin liian aikaisin, minulla oli vapaata aikaa katsella
matkustajien liikettä asemalla. Jotkut jättivät tai noutivat käsimatkatavaransa
\VUgras{(matkatavara)säilytyksessä} \textsuperscript{(25)}; toiset tutkivat
\VUgras{aikatauluja} \textsuperscript{(26)}. Saapuvat matkustajat kiiruhtivat
\VUgras{uloskäynnille} \textsuperscript{(32)} \VUgras{matkalaukkunsa}
\textsuperscript{(24)} kädessään. Jotkut odottivat
\VUgras{(matkatavara)luovutuksessa} \textsuperscript{(27)} ja näyttivät
\VUgras{merkkiään} \textsuperscript{(29)} saadakseen takaisin matka-arkkunsa,
\VUgras{laatikkonsa} \textsuperscript{(28)} tai \VUgras{korinsa}
\textsuperscript{(30)}.

%%K t12-07-1 p
7. --- \VUgras{Aksiisivirkailijat} \textsuperscript{(33)} tutkivat huolellisesti
pakkaukset todetakseen, onko jotain tullattavaa.

%%K t12-08-1 p
8. --- Jotkut matkustajat menivät \VUgras{ravintolaan} \textsuperscript{(34)},
jossa \VUgras{tarjoilija} \textsuperscript{(35)} ponnisteli palvellakseen heitä.
Heidän täytyi kiirehtiä, sillä \VUgras{juna} \textsuperscript{(36)}, joka on
pikajuna, ei pysähdy kauaksi aikaa asemalle.

%%K t12-09-1 p
9. --- Sitten menin lähtölaiturille ja etsin läpikulkevaa \VUgras{vaunua}
\textsuperscript{(37)}, jotta minun ei tarvitsisi vaihtaa vaunua matkan aikana.
Nousin käytävävaunuun, jossa on \VUgras{osasto} \textsuperscript{(38)}
tupakoiville ja osasto, joka on varattu « yksinäisille naisille ».

%%K t12-tit-2 sub
\VUpk{10.2pt}{40}{Lähtö yöllä.}
\VUsaut{3.0mm}

%%K t12-10-1 p
10. --- Noustuani \VUgras{astinlaudalle} \textsuperscript{(40)}, suljettuani
\VUgras{oven} \textsuperscript{(39)}, kierrettyäni ylös \VUgras{verhon}
\textsuperscript{(41)} ja pantuani käsimatkatavarani \VUgras{verkkoon}
\textsuperscript{(43)}, istuuduin \VUgras{penkille} \textsuperscript{(42)}.
Kun näin \VUgras{lampunhoitajan} \textsuperscript{(44)} sytyttävän lamput,
ajattelin, että minun täytyisi viettää yö vaunussa, ja kutsuin sen virkailijan,
joka vastaa \VUgras{päätyynyjen} \textsuperscript{(45)} vuokrauksesta,
pyytääkseni yhtä niistä.

%%K t12-11-1 p
11. --- Junan konduktööri tuli tarkastamaan liput. Kysyin häneltä, miksi emme
vielä olleet lähteneet, sillä on oikea aika. Hän vastasi minulle, että
\VUgras{junapäällikkö} \textsuperscript{(46)} on panettamassa polkupyöriä
\VUgras{matkatavaravaunuun} \textsuperscript{(47)}. Sitä paitsi kaikkia
jalkalämmittimiä \textsuperscript{(48)} ei ollut vielä vaihdettu.

%%K t12-12-1 p
12. --- Mutta lähtisimme pian, sillä \VUgras{lämmittäjä} \textsuperscript{(50)}
otti \VUgras{tenderillään} \textsuperscript{(49)} hiiltä yllyttääkseen
\VUgras{veturin} \textsuperscript{(54)} tulta, ja \VUgras{koneenkäyttäjä}
\textsuperscript{(51)} odotti vain merkkiä \VUgras{aliasemapäälliköltä}
\textsuperscript{(52)}, joka seisoi \VUgras{laiturilla} \textsuperscript{(53)}.

%%K t12-13-1 p
13. --- Veturin \VUgras{höyrykattila} \textsuperscript{(56)} jyrisi kumeasti;
\VUgras{savupiippu} \textsuperscript{(55)} sylki virtoja savua sekoittuneena
\VUgras{höyryyn} \textsuperscript{(60)}. Kimeä \VUgras{vihellys}
\textsuperscript{(59)} kaikui, ja \VUgras{männät} \textsuperscript{(58)}
alkoivat liikkua ja pyörittivät raskaan koneen \VUgras{pyöriä}
\textsuperscript{(57)}.

%%K t12-tit-3 sub
\VUsaut{3.0mm}
\VUpk{10.2pt}{40}{Matkaan!}
\VUsaut{3.0mm}

%%K t12-14-1 p
14. --- Junan vauhti, sen kulkiessa \VUgras{kiskoilla} \textsuperscript{(64)},
kiihtyi; \VUgras{laiturit} \textsuperscript{(65)} katosivat; ohitimme
\VUgras{käymälät} \textsuperscript{(66)}, ja myös rivin vaunuja, jotka olivat
asetettuina toiselle \VUgras{raiteelle} \textsuperscript{(63)}:
\VUgras{makuuvaunuja} \textsuperscript{(67)}, \VUgras{ravintolavaunun}
\textsuperscript{(68)}, \VUgras{postivaunun} \textsuperscript{(69)}.

%%K t12-noto-1 noto
\VUnotes{143.2mm}{%
\textit{(*) Huom.} --- \textit{Itsestään selvää on, että matkan kuvaus noudattaa
sotaa edeltänyttä rajaa.}%
}

%%K t12-15-1 p
15. --- Sitten \VUgras{tavara-asema} \textsuperscript{(71)} liukui ohi silmiemme.
Vajojen alla virkailijat lastasivat \VUgras{tavaroita} \textsuperscript{(72)},
jotka lähetetään tavarajunalla. \VUgras{Vaihtomiehet} \textsuperscript{(76)}
\VUgras{vesisäiliön} \textsuperscript{(73)} luona vaihtoivat
\VUgras{karjavaunuja} \textsuperscript{(75)} ja \VUgras{lavavaunuja}
\textsuperscript{(79)} muodostaakseen \VUgras{tavarajunan}
\textsuperscript{(80)}.

%%K t12-16-1 p
16. --- Ohitimme viimeisen \VUgras{vaihteen} \textsuperscript{(78)}, jonka luona
vaihdemies seisoi; kuljimme \VUgras{opastinlevyn} \textsuperscript{(86)} ohi, ja
asema katosi kaukaisuuteen.

%%K t12-17-1 p
17. --- Pian kuljimme \VUgras{rautasillan} \textsuperscript{(74)} yli, joka
kulkee \VUgras{joen} \textsuperscript{(81)} yli. Kuljettuamme tuon sillan yli
seurasimme virtaa, jolla voimakkaat \VUgras{hinaajat} \textsuperscript{(82)}
vetivät raskaita \VUgras{proomuja} \textsuperscript{(83)}. Näimme lisäksi
\VUgras{matkustaja-aluksen} \textsuperscript{(84)}, kun se laski
\VUgras{ponttoniin} \textsuperscript{(85)}.

%%K t12-18-1 p
18. --- Ohitettuamme raivoisan nopeasti useita asemia kuljimme muutamien
\VUgras{tunnelien} \textsuperscript{(87)} läpi ja jätimme taaksemme
lukemattomia \VUgras{pieniä taloja} \textsuperscript{(88)} radanvartijoille.
Junan tärinän tuudittamana nukahdin syvään.

%%K t12-tit-4 sub
\VUsaut{3.0mm}
\VUpk{10.2pt}{40}{Deutsch-Avricourtin asema.}
\VUsaut{3.0mm}

%%K t12-19-1 p
19. --- Minut herätti äkkiä virkailijan ääni, joka huusi:
« Deutsch-Avricourt! \textsuperscript{(*)}. Kaikki pois! » Tullin tarkastus ja
kaikki rajan väsyttävät muodollisuudet tapahtuivat. Minun täytyi asettaa
taskukelloni aseman \VUgras{suuren kellon} \textsuperscript{(89)} mukaan, joka
kävi viisikymmentäviisi minuuttia edellä; tuskin hereillä erotin vaivoin
\VUgras{suuren viisarin} \textsuperscript{(90)} \VUgras{pienestä}
\textsuperscript{(91)}; itse \VUgras{kellotaulu} \textsuperscript{(92)} oli
aivan sekaisin aamusumusta.

%%K t12-20-1 p
20. --- Samalla kun tullivirkailijat tekivät tarkastustaan, tulkitsin haukotellen
niitä \VUgras{julisteita} \textsuperscript{(93)}, jotka koristavat aseman
muureja. Söin aamiaisen kuluttaakseni aikaa, tutkin karttaa saksalaisesta
\VUgras{verkosta} \textsuperscript{(94)} nähdäkseni, kuinka pitkä matka minut
vielä erottaa Strasbourgista, ja menin takaisin vaunuuni vahvistuneena
toivosta saavuttaa pian matkani päämäärä.
\VUsaut{6.0mm}
\VUfilet{22mm}
